\section{Implementasi Hash Table di C}
\label{sec:ch15-hash}

\subsection{Mengapa Hash Table}

Pencarian linear pada daftar simbol memiliki kompleksitas
\(O(n)\) per lookup. Untuk program dengan ribuan identifier, biaya
ini menjadi signifikan. Hash table menurunkan kompleksitas rata-rata
menjadi \(O(1)\) jika fungsi hash dan ukuran tabel dipilih dengan
baik.

\subsection{Fungsi Hash}

Fungsi hash yang sederhana untuk string:
\[
h(s) = \left(\sum_{i=0}^{n-1} s_i \cdot 31^{n-1-i}\right)
\bmod m,
\]
dengan \(m\) ukuran tabel. Perkalian dengan 31 menyebar karakter
dengan cukup baik dan mudah diimplementasikan sebagai
\code{h = h * 31 + c}.

\subsection{Chaining}

Setiap bucket berisi daftar simbol dengan hash yang sama. Operasi
\code{define} menambahkan ke depan daftar; \code{lookup} menelusuri
daftar bucket. Ukuran tabel sebaiknya bilangan prima atau pangkat
dua dengan fungsi hash yang sesuai.

\begin{figure}[H]
\centering
\begin{tikzpicture}[node distance=4mm]
  \node[memorycell, minimum width=2.2cm] (b0) {bucket 0};
  \node[memorycell, minimum width=2.2cm, below=of b0] (b1) {bucket 1};
  \node[memorycell, minimum width=2.2cm, below=of b1] (b2) {bucket 2};
  \node[phase, minimum width=2.2cm, right=of b0] (s0) {\code{x}};
  \node[phase, minimum width=2.2cm, right=of b1] (s1) {\code{y}};
  \node[phase, minimum width=2.2cm, right=of s1] (s2) {\code{z}};
  \draw[arrow] (b1) -- (s1);
  \draw[arrow] (s1) -- (s2);
\end{tikzpicture}
\caption{Hash table dengan chaining}
\label{fig:ch15-hash}
\end{figure}

\subsection{Implementasi Proyek}

Pada proyek Subset-C, tabel simbol diimplementasikan sebagai daftar
tertaut per scope. Struktur ini mudah dipahami dan cukup untuk
program kecil. Mahasiswa dapat menggantinya dengan hash table sebagai
pengayaan.

Antarmuka tabel simbol ditunjukkan pada \lstref{lst:ch15-symtab-h}.

\includecode[
  language=C,
  firstline=1,
  lastline=45,
  firstnumber=1,
  caption={Antarmuka \file{symtab.h}},
  label={lst:ch15-symtab-h}
]{../../code/subsetc/include/symtab.h}

Fungsi \code{symtab\_define} dan \code{symtab\_lookup} menjadi inti
implementasi. Keduanya menangani nama yang tidak selalu
null-terminated dengan menerima panjang nama secara eksplisit.
